\chapter[Sermon 51. The Thanksgiving of the Elders Is Expounded, the Temple Is Opened in Heaven, the Ark Appears, and There Were Made Lightnings, Etc.]{The Thanksgiving of the Elders Is Expounded, the Temple Is Opened in Heaven, the Ark Appears, and There Were Made Lightnings, Etc.}
\label{sermon:051}
\markright{Sermon 51. The Thanksgiving of the Elders Is Expounded, the ...}

And the Gentiles were angry, and your wrath has come, and the time for the dead to be judged, so that you may give reward to your servants the prophets and saints, and to those who fear your name, both small and great; and to destroy those who destroy the earth. And the temple of God was opened in heaven, and there was seen in his temple the Ark of his Testament. And there followed lightnings, and voices, and thunderings, and an earthquake, and a great hail.

I showed you how the Elders gave thanks to God for their salvation, also praising God’s righteousness and excellent truth, which he demonstrates in his judgment as most righteous, wherein he rewards the godly with just rewards and punishes the wicked with deserved punishments. And by this figure of speech, they teach us that judgment will assuredly come, and that it will be entirely holy and just. Would that those who judge them to be speaking of trifles would diligently consider these things themselves, those who speak of that horrible and most dreadful day of judgment. For we look for things more terrible than any tongue, however eloquent, is able to express.

He recounts the wrath or tyranny of non-believers, against the faithful, cruelly and continually executed, so truly that to many it seemed God was unresponsive, and neither could nor would be angry. But the judgment ones made, the elders extol God’s truth, and say the wrath has come. Doubtless the holy Prophets of God have always threatened punishments, testifying that God is angry, both with sinners and with sins. But where the wrath of God did not appear immediately, the Prophets seemed to frighten people with empty terrors, as if making them afraid of their shadows. But now, say the elders, the truth has appeared, and the wrath of God is come. And the wrath of God shows itself in God’s just vengeance.

\index{Resurrection}Moreover, they also praise this for demonstrating God’s truth and justice, for the time of the dead has arrived, that they be judged. Until now, while the world prospered, they seemed to tell fables and old women’s tales, which spoke of the resurrection of the dead and the life to come. For the resurrection of the dead was scorned by philosophers and worldly men. But the elders rejoice also, that this time has come, and that the dead are raised up, that is, that the bodies of the dead have risen again and come to judgment. Concerning this, the Apostle says: “We must all appear openly before the judgment seat of God, that every one may receive according to what the body has done, whether good or evil.” 2 Corinthians, chapter 5.

Furthermore, they greatly praise God’s justice and truthfulness when they clearly explain how God, by his just judgment, renders to every one according to his deeds. He declares what he rewards and whom he rewards. First, he pays wages or hire. For reward is promised by God for good works. For in Jeremiah 13, the Lord says, “Restrain your voice from weeping, for there is a reward for your work.” And the Lord also says in the gospel, “Be glad and rejoice, for your reward is great in heaven.” And again, “The Son of Man shall come in the glory of his Father with his angels, and then shall he render to every one according to his doings.” So the Apostle said that every one must rise in his own body, that every one may receive such things as are done by the body, whether it be good or evil. While this world flourishes, and the wicked rejoice in their indulgence, and the godly are afflicted, and afflict themselves with continual mortification, the flesh judges that these lose both labor and expense, while the others are very happy. This too is declared in chapters 3 and 4 of Malachi. But at the final judgment it shall finally appear that the godly have not labored in vain, nor have the wicked scorned God unpunished and despised godliness. For God rewards every one according to the quality of his work, which he calls wages. Nevertheless, the godly do not abuse this saying in the meantime, acknowledging it to be of free mercy that they have believed and worked with good faith, and that their good work is therefore accepted by God because it is in Christ. Of this I have written in book 3, chapter 10, concerning the grace of God justifying, showing that reward cannot be proven by merit.

Secondly, they declare to whom he gives reward; I say to two sorts of men: good, I mean, and evil. Again, he recounts many kinds of good men. First, he calls these the servants of God, as those who are subject to the empire of God alone and obey him in all things. Immediately, he names the prophets, teachers of churches. Concerning their state, more is spoken in chapter 11. These are sometimes more unfortunate than any others in this world and are accounted by many as great offenders, so that if they were removed, all clarity would seem to return. Therefore, they are justly recounted in the register of those who receive a reward from the Lord, namely, in recompense for their labor. Now into this account come also the saints, that is to say, all the godly who, being sanctified through faith with the spirit and blood of God, have lived a holy life, keeping themselves from all worldly pollution. Moreover, into the reward and honor of the holy saints are reckoned those who fear the name of the Lord: that is, they who are very holy and religious in deed. Finally, lest any man should think any of the faithful excluded, he adds, “to small and great”: that is to say, to men of all ages, status, and sex, and so forth.

After he comes to the evil, and adds: “and shouldest destroy those who destroyed the earth.” These things seem borrowed from the prophets, with whom there is much mention of those who destroy the earth, whom the Lord will ultimately destroy. And under the name of destroyers, John understands first Tyrants, Kings, and Princes, who are persecutors of the church. Also, men of war and soldiers, who by unjust wars destroy all things with sword and fire. Secondly, he understands unjust judges, moreover oppressors of the poor, who afflict widows and orphans; moreover, those who through usury, theft, deceitfulness, extortion, and evil means are harmful to all men, and by their insatiable covetousness bring about a scarcity of all things. Finally, those who by whoredom and adultery defile and break holy matrimony. Last, heretics destroy the earth, and those who infect people with corrupt doctrine, who dwell upon the earth; into this number also come seditious persons and traitors, and other wicked men.

These the Lord will destroy with everlasting destruction, so that they do not cease to be those who perish, but become much more miserable while they are tormented with torments that will never end. Those who are wasteful and extravagant are said to be lost, yet in perishing thus they do not cease to be, but daily proceed to be more miserable, which is destruction itself.

\index{John, Saint}Furthermore, Saint John questions this doctrine of the reward of the godly; and what he previously presented as thankful praise and a joyous triumph, he now proposes as if it were to be seen with the eyes through a heavenly vision. He concludes this vision with the opening of the Temple, just as he began it with the opening of Heaven. For the loving Lord opens to his servants heaven itself to be seen by the mind, so that we should doubt nowhere the glory prepared for us in Heaven; neither should say, “Who has seen those heavenly things that are promised us?” For just as the blessed fathers, the Prophets and Apostles have had many visions of this sort, effective, true, and godly, so may each of us with the eyes of our mind, through true faith, look into Heaven itself. I know well that worldly men pass nothing upon such visions, as the Lord said in the Gospel: “The world cannot receive the Spirit of truth, for that it sees him not, neither knows him.” Let us not care for their contempt.

Let us therefore, see what is prepared for the servants of God in another world. First, John saw heaven open: and in heaven itself, he sees also the very temple of God, open, that is, to all the godly. By the Temple of God, he understands the secrets of God, the inward and privy parts of Heaven, whereinto he will receive to the fruition of himself all believers. But in that divine temple of heaven was seen the Ark of his Testament. For God made a covenant or league with the faithful, that he would be their God, their fullness, and a most plentiful sea of all goodness, a most abundant, and most sufficient plenty of all things. The confirmation, testimony, and declaration whereof is the Ark of covenant, the very Son of God, in whom dwells all fullness of deity, and in whom we are made perfect. For he is the Ark, in whom are laid up all celestial treasures, full of grace and truth. This Ark of good things, and of eternal felicity, appears in heaven. For the Son of God is on the throne of God. The liberal and bountiful heavenly Father will pour out this Ark upon his children, granting to them through Christ his only Son all heavenly gifts, that we might be partakers of all Christ’s benefits, even to the deity, wherein he excels his brethren. Hereby it appears how Moses prepared the Ark, after the example of the same which he saw in heaven: and the figure whereof was the Ark of the covenant, \&c. Otherwise we shall hear in the twenty-first chapter of this book, that there is no temple in heaven, \&c.

These most beautiful things to be seen, and most pleasant to be heard, the Son of God has set forth to be seen and heard by us. Consequently, he adds that punishments are prepared for the wicked, and explains them in various ways, presenting them to be seen. Hitherto, there had been in the world lightnings, voices, and thunderings, etc. The Holy Spirit shone to the world, drawing through the doctrine of the truth, moving and inspiring fear. But the reckless world would not understand, nor even hear the manner and way of salvation; therefore, divine justice requires that they be addressed in another language, and therefore, by God’s just judgment, are now made lightnings, etc. And by this heap of words, he signifies the horrible punishment that God will inflict on the wicked. He appears to have alluded to the burning of Sodom, also to the words of the godly Prophet: “It shall reign upon sinners snares of fire, brimstone, and spirit of tempest,” in Psalm 11. Therefore, this vision is concluded, as in the story of Saint Matthew’s Gospel: and these shall go into everlasting punishment, and the righteous into life everlasting.

In these last eight chapters, we have the third part of this book, and a significant summary of the church’s history from the time of John until the end of the world. Through this, we are instructed in the true faith and warned of all dangers and betrayals by which the true faith is attacked, so that, being watchful, we may avoid all corruption and cunning deception and be kept safe. To God be praise and glory.
